% GENERATED FILE. Edit canonical lessons, then run npm run generate:books.
% canonical-source-hash: fnv1a64:ce7629d172b48d3d
% canonical-lessons: TA-C461-porumai, TA-C461-nitanam, TA-C461-caki, TA-C461-arvam, TA-C461-aval

\chapter{Patience and Eagerness}
\label{ch:ta-patience-and-eagerness}
\hlchaptermodality{461}

\begin{chapteropening}
\textbf{By the end of this chapter:} \emph{I can talk about patience, calm, interest and eagerness.}

Complete the last lesson of chapter 461: I can talk about patience, calm, interest and eagerness.
\end{chapteropening}

\section[\texorpdfstring{\ta{பொறுமை} (poṟumai)}{பொறுமை (poṟumai)}]{\ta{பொறுமை} (poṟumai) — patience}

\label{lesson:TA-C461-porumai}



\begin{quote}
Before the new one: say the Tamil for ecstasy, then the Tamil for to tingle.
\end{quote}

\subsection*{You'll want to know: \ta{பொறுமை}}
\textbf{\ta{பொறுமை}} — \emph{poṟumai} — \textquotedblleft{}patience\textquotedblright{}.

\begin{quote}
\ta{குழந்தைகளிடம்} \ta{ஆசிரியருக்குப்} \ta{பொறுமை} \ta{வேண்டும்}.
\end{quote}

\begin{quote}
\emph{kuḻantaikaḷiṭam āciriyarukkup poṟumai vēṇṭum.}
\end{quote}

\begin{quote}
A teacher needs patience with children.
\end{quote}

\ta{பொறுமையாக} \ta{இரு}, to be patient.

\begin{cousinweb}
\textbf{\ta{பொறுமை}} (\emph{poṟumai}) is native Dravidian, a noun from \ta{பொறு} \emph{poṟu}, \textquotedblleft{}to bear, to endure, to be patient\textquotedblright{} (DEDR 4565).
\end{cousinweb}

\begin{grammarlens}[title={What you've built}]
One more word for describing what you lived and felt.
\end{grammarlens}

\subsection*{Your turn}
\begin{itemize}
\raggedright
  \item \emph{Say it:} \emph{poṟumai}
  \item \emph{Say it:} \emph{poṟumai}, once more
  \item \emph{Say it:} \emph{kuḻantaikaḷiṭam āciriyarukkup poṟumai vēṇṭum.} — A teacher needs patience with children.
  \item \emph{Recall:} say the Tamil for ecstasy, then the Tamil for to tingle, then say \emph{poṟumai} again
\end{itemize}

\begin{tcolorbox}[breakable,colback=teal!4,colframe=teal!35!black,title={Before you move on}]
What does \ta{பொறுமை} mean? (\textquotedblleft{}Patience\textquotedblright{}.) Say it once more.
\end{tcolorbox}
\section[\texorpdfstring{\ta{நிதானம்} (nitāṉam)}{நிதானம் (nitāṉam)}]{\ta{நிதானம்} (nitāṉam) — composure; steadiness}

\label{lesson:TA-C461-nitanam}



\begin{quote}
Before the new one: say the Tamil for to tingle, then the Tamil for patience.
\end{quote}

\subsection*{You'll want to know: \ta{நிதானம்}}
\textbf{\ta{நிதானம்}} — \emph{nitāṉam} — \textquotedblleft{}composure; steadiness\textquotedblright{}.

\begin{quote}
\ta{அவர்} \ta{கோபப்படாமல்} \ta{நிதானமாகப்} \ta{பதில்} \ta{சொன்னார்}.
\end{quote}

\begin{quote}
\emph{avar kōpappaṭāmal nitāṉamākap patil coṉṉār.}
\end{quote}

\begin{quote}
He answered calmly, without getting angry.
\end{quote}

\ta{நிதானமாக}, calmly; \ta{நிதானம்} \ta{இழந்தார்}, he lost his composure.

\begin{cousinweb}
\textbf{\ta{நிதானம்}} (\emph{nitāṉam}) is borrowed from Sanskrit \emph{nidāna}, \textquotedblleft{}a band, a halter; a first cause\textquotedblright{} (CDIAL 7196); Tamil gives it the sense of measured calm.
\end{cousinweb}

\begin{grammarlens}[title={What you've built}]
One more word for describing what you lived and felt.
\end{grammarlens}

\subsection*{Your turn}
\begin{itemize}
\raggedright
  \item \emph{Say it:} \emph{nitāṉam}
  \item \emph{Say it:} \emph{nitāṉam}, once more
  \item \emph{Say it:} \emph{avar kōpappaṭāmal nitāṉamākap patil coṉṉār.} — He answered calmly, without getting angry.
  \item \emph{Recall:} say the Tamil for to tingle, then the Tamil for patience, then say \emph{nitāṉam} again
\end{itemize}

\begin{tcolorbox}[breakable,colback=teal!4,colframe=teal!35!black,title={Before you move on}]
What does \ta{நிதானம்} mean? (\textquotedblleft{}Composure; steadiness\textquotedblright{}.) Say it once more.
\end{tcolorbox}
\section[\texorpdfstring{\ta{சகி} (caki)}{சகி (caki)}]{\ta{சகி} (caki) — to tolerate; to put up with}

\label{lesson:TA-C461-caki}



\begin{quote}
Before the new one: say the Tamil for patience, then the Tamil for composure.
\end{quote}

\subsection*{You'll want to know: \ta{சகி}}
\textbf{\ta{சகி}} — \emph{caki} — \textquotedblleft{}to tolerate; to put up with\textquotedblright{}.

\begin{quote}
\ta{அவனுடைய} \ta{கெட்ட} \ta{பழக்கத்தை} \ta{என்னால்} \ta{சகிக்க} \ta{முடியவில்லை}.
\end{quote}

\begin{quote}
\emph{avaṉuṭaiya keṭṭa paḻakkattai eṉṉāl cakikka muṭiyavillai.}
\end{quote}

\begin{quote}
I could not tolerate his bad habit.
\end{quote}

Past (I …): \ta{சகித்தேன்}, \emph{cakittēṉ}.

With \ta{என்னால்} … \ta{முடியவில்லை}: I could not bear it.

\begin{cousinweb}
\textbf{\ta{சகி}} (\emph{caki}) is borrowed from the Sanskrit verb \emph{sahatē}, \textquotedblleft{}bears, suffers\textquotedblright{} (CDIAL 13304); Tamil turns the root into a verb in -\ta{இ}.
\end{cousinweb}

\begin{grammarlens}[title={What you've built}]
One more word for describing what you lived and felt.
\end{grammarlens}

\subsection*{Your turn}
\begin{itemize}
\raggedright
  \item \emph{Say it:} \emph{caki}
  \item \emph{Say it:} \emph{caki}, once more
  \item \emph{Say it:} \emph{avaṉuṭaiya keṭṭa paḻakkattai eṉṉāl cakikka muṭiyavillai.} — I could not tolerate his bad habit.
  \item \emph{Recall:} say the Tamil for patience, then the Tamil for composure, then say \emph{caki} again
\end{itemize}

\begin{tcolorbox}[breakable,colback=teal!4,colframe=teal!35!black,title={Before you move on}]
What does \ta{சகி} mean? (\textquotedblleft{}To tolerate; to put up with\textquotedblright{}.) Say it once more.
\end{tcolorbox}
\section[\texorpdfstring{\ta{ஆர்வம்} (ārvam)}{ஆர்வம் (ārvam)}]{\ta{ஆர்வம்} (ārvam) — keenness; interest}

\label{lesson:TA-C461-arvam}



\begin{quote}
Before the new one: say the Tamil for composure, then the Tamil for to tolerate.
\end{quote}

\subsection*{You'll want to know: \ta{ஆர்வம்}}
\textbf{\ta{ஆர்வம்}} — \emph{ārvam} — \textquotedblleft{}keenness; interest\textquotedblright{}.

\begin{quote}
\ta{என்} \ta{மகளுக்கு} \ta{இசையில்} \ta{நிறைய} \ta{ஆர்வம்} \ta{இருக்கிறது}.
\end{quote}

\begin{quote}
\emph{eṉ makaḷukku icaiyil niṟaiya ārvam irukkiṟatu.}
\end{quote}

\begin{quote}
My daughter is very keen on music.
\end{quote}

The thing one is keen on takes -\ta{இல்}: \ta{இசையில்} \ta{ஆர்வம்}.

\begin{cousinweb}
\textbf{\ta{ஆர்வம்}} (\emph{ārvam}) is native Dravidian (DEDR 381), where it means affection and desire.
\end{cousinweb}

\begin{grammarlens}[title={What you've built}]
One more word for describing what you lived and felt.
\end{grammarlens}

\subsection*{Your turn}
\begin{itemize}
\raggedright
  \item \emph{Say it:} \emph{ārvam}
  \item \emph{Say it:} \emph{ārvam}, once more
  \item \emph{Say it:} \emph{eṉ makaḷukku icaiyil niṟaiya ārvam irukkiṟatu.} — My daughter is very keen on music.
  \item \emph{Recall:} say the Tamil for composure, then the Tamil for to tolerate, then say \emph{ārvam} again
\end{itemize}

\begin{tcolorbox}[breakable,colback=teal!4,colframe=teal!35!black,title={Before you move on}]
What does \ta{ஆர்வம்} mean? (\textquotedblleft{}Keenness; interest\textquotedblright{}.) Say it once more.
\end{tcolorbox}
\section[\texorpdfstring{\ta{ஆவல்} (āval)}{ஆவல் (āval)}]{\ta{ஆவல்} (āval) — eagerness; longing}

\label{lesson:TA-C461-aval}



\begin{quote}
Before the new one: say the Tamil for to tolerate, then the Tamil for keenness.
\end{quote}

\subsection*{You'll want to know: \ta{ஆவல்}}
\textbf{\ta{ஆவல்}} — \emph{āval} — \textquotedblleft{}eagerness; longing\textquotedblright{}.

\begin{quote}
\ta{உன்னைப்} \ta{பார்க்க} \ta{நான்} \ta{ஆவலாக} \ta{இருக்கிறேன்}.
\end{quote}

\begin{quote}
\emph{uṉṉaip pārkka nāṉ āvalāka irukkiṟēṉ.}
\end{quote}

\begin{quote}
I am eager to see you.
\end{quote}

With an infinitive: \ta{பார்க்க} \ta{ஆவல்}, eager to see.

\begin{cousinweb}
\textbf{\ta{ஆவல்}} (\emph{āval}) is native Dravidian, a noun from \ta{ஆவு} \emph{āvu}, \textquotedblleft{}to desire\textquotedblright{} (DEDR 394).
\end{cousinweb}

\begin{grammarlens}[title={What you've built}]
One more word for describing what you lived and felt.
\end{grammarlens}

\subsection*{Your turn}
\begin{itemize}
\raggedright
  \item \emph{Say it:} \emph{āval}
  \item \emph{Say it:} \emph{āval}, once more
  \item \emph{Say it:} \emph{uṉṉaip pārkka nāṉ āvalāka irukkiṟēṉ.} — I am eager to see you.
  \item \emph{Recall:} say the Tamil for to tolerate, then the Tamil for keenness, then say \emph{āval} again
\end{itemize}

\begin{tcolorbox}[breakable,colback=teal!4,colframe=teal!35!black,title={Before you move on}]
What does \ta{ஆவல்} mean? (\textquotedblleft{}Eagerness; longing\textquotedblright{}.) Say it once more.
\end{tcolorbox}
